\section{Conclusion}
\label{sec:conclusion}

PSBD transfers to ViTs once its perturbation is placed where a ViT carries a backdoor. Our reading of the original placement, dropout on the residual stream, detects the attacks with a firm shortcut and misses the rest. Moving the same dropout to the attention input already repairs much of that, and masking whole tokens there, which we call PSBD-TM, reaches a mean AUROC of 0.935 on ViT-B/16 and 0.968 on Swin-S, with the largest gains at 1\% poisoning, where a detector is needed most.

The placement is not a tuning accident. A backdoor in a ViT is 1 direction in the residual stream, written in the last third of the network and routed through attention from the trigger's patch tokens to the class token. A perturbation that removes tokens before attention removes the shortcut, while a perturbation on the residual stream perturbs the decision itself. The same account explains why noise and masking swap places between the attention input and the MLP input, predicts which attacks show the prediction shift toward the target class that the original paper observed, and reproduces the direction removal of Karayal\c{c}{\i}n et al.\ once the direction is removed from every residual write.

For a practitioner the recommendation is short. On a ViT or a Swin, run PSBD with token masking at the attention input and choose the rate by the clean shift ratio as the original does. Against a knowledgeable attacker, combine several placements, since an attacker who evades 1 probe does not evade the others. The union of probes costs a sweep per probe and restores detection above 0.97 on every attack we trained against a single probe.

3 questions remain open. An attacker who trains against the union of probes, a like-for-like comparison to other input-level detectors beyond the smoke test, and the all-to-all attack, whose backdoor has no single direction and which, on the present evidence, needs a different statistic.
